\section{TensorRT}

\subsection{Kiến trúc TensorRT}

TensorRT là bộ SDK tối ưu và thực thi suy luận cho GPU NVIDIA. Builder tiếp nhận
đồ thị mạng, áp dụng tối ưu như fusion, lựa chọn kernel và precision, sau đó tạo
serialized engine. Runtime nạp engine, cấp phát tensor và thực thi qua execution
context \cite{tensorrt}.

Engine phụ thuộc vào kiến trúc GPU, phiên bản TensorRT, plugin và shape profile.
Vì vậy, manifest của model artifact phải ghi rõ môi trường tạo engine và điều kiện
tương thích.

\subsection{ONNX}

ONNX là định dạng trung gian giúp trao đổi mô hình giữa framework huấn luyện và
runtime suy luận \cite{onnx}. Pipeline của đồ án export checkpoint sang ONNX,
kiểm tra graph, so sánh đầu ra với framework gốc rồi mới xây dựng TensorRT
engine. Các toán tử không được hỗ trợ phải được thay thế, phân rã hoặc hiện thực
bằng plugin.

\subsection{TensorRT Engine}

TensorRT engine là artifact đã được tối ưu cho một tập cấu hình cụ thể. Quá trình
build có thể lựa chọn tactic khác nhau theo workspace và phần cứng. Engine cần
đi kèm input shape, data type, tên binding, preprocessing, postprocessing và hàm
băm để Edge Runtime tải đúng cách.

\subsection{FP16 Inference}

FP16 inference cho phép TensorRT dùng kernel bán chính xác ở các tầng được hỗ
trợ và quay về precision khác khi cần. Đây thường là cấu hình triển khai đầu tiên
vì quy trình chuyển đổi đơn giản hơn INT8 và mức suy giảm accuracy nhỏ. Kết quả
vẫn phải được so sánh với FP32 trên cùng tập test.

\subsection{INT8 Inference}

INT8 inference yêu cầu scale activation phù hợp. Với PTQ, TensorRT dùng
calibrator; với QAT, thông tin lượng tử được mang theo từ graph nếu pipeline
export hỗ trợ. Một số tầng có thể giữ FP16 hoặc FP32 để bảo toàn accuracy, tạo ra
engine mixed precision.

\subsection{TensorRT trên NVIDIA Jetson}

Trên Jetson, TensorRT được phân phối cùng hệ sinh thái JetPack và sử dụng CUDA,
cuDNN cùng driver tương ứng. Việc build engine trực tiếp trên thiết bị đích giúp
giảm rủi ro không tương thích. Đồ án lưu riêng engine cho Jetson Nano và Jetson
Orin Nano, thay vì giả định một engine có thể dùng chung cho mọi GPU.